\section{Related work}
\label{sec:related}

\paragraph{Team formation and composition.} Choosing a group of agents whose joint performance exceeds what their individual abilities suggest is studied in operations research, multiagent systems and organizational psychology. \citet{boon2003} formulate team formation in sport as matching the quality a team requires at each position with the quality candidate players supply, and \citet{pantuso2017} treats football squad composition as a stochastic program over uncertain future player values subject to budget and regulatory constraints. In multiagent systems, \citet{andrejczuk2016} and \citet{andrejczuk2018} review team composition and formation, distinguishing problems with and without explicit roles and connecting computational models to evidence from organizational psychology. Our problem has explicit roles (positional quotas) and an explicit performance model, and differs from most of this literature in that the value of a team is predicted by a model learned from game outcomes rather than specified as a known function of individual attributes.

\paragraph{Player value and lineup interaction.} Adjusted plus-minus regresses scoring margin on indicators for the players on the floor \citep{rosenbaum2004}, and ridge regularization makes the estimates usable out of sample \citep{sill2010}; \citet{macdonald2011} adapts the approach to hockey. These models are additive in players. Several authors have tried to measure the non-additive part. \citet{maymin2013} simulate lineups from player skill ratings and compare lineup performance with the sum of its parts, and \citet{pelechrinis2018} embeds a network of lineups to predict which of two lineups will outscore the other. Interaction effects are much harder to estimate than individual effects, and we treat the existence of usable interaction signal in our model as an empirical question (Section~\ref{sec:validation}) rather than an assumption.

\paragraph{Graph neural networks and adversarial evaluation.} Graph convolutional networks \citep{kipf2017} and their relational variants with separate weights for different edge types \citep{schlichtkrull2018} propagate information along observed relationships, and DeepSet pooling \citep{zaheer2017} produces representations of sets that do not depend on the order of their elements. Graph models of players and their spatial relationships have been used for tactical analysis in football \citep{wang2024tacticai}. Leader-follower (Stackelberg) games, in which one player commits to a strategy and the other best-responds, are standard models for planning against an adaptive opponent \citep{conitzer2006,tambe2011}. We use the leader-follower structure for the deployment problem only. Roster selection itself is treated as a search over a fixed scoring function.
